\section{Method}
\label{sec:method}

We present AnomalyClaw, a training-free framework for cross-domain visual anomaly detection built on three components: an \emph{extensible expert pool} that provides quantitative grounding, an \emph{adversarial debate engine} that enables self-verification through structured argumentation, and an \emph{autonomous controller} that orchestrates expert selection and debate depth.

\subsection{Problem Formulation}
\label{sec:formulation}

Given a query image $\query$, a domain code $\domcode$ identifying the inspection context (e.g., \texttt{mvtec\_metal\_nut}, \texttt{brain\_mri}), and a bank of normal images $\bank_\domcode = \{I^n_1, \ldots, I^n_N\}$, the goal is to produce an anomaly verdict $v \in \{\texttt{normal}, \texttt{anomalous}\}$ and a set of anomaly claims with confidence scores.
The system operates in a \emph{training-free, few-shot} setting: no parameters are learned from anomaly examples, and only $k$ normal references (typically $k \leq 5$) are used per query at inference time.

\subsection{Extensible Expert Pool}
\label{sec:experts}

The expert pool $\Phi = \{\expert_1, \ldots, \expert_E\}$ is a collection of lightweight analysis modules, each producing structured text evidence about the query image.
Any expert that takes an image (and optionally reference images) as input and produces a textual report can be integrated without modifying the debate engine.

\paragraph{Patch-level expert.}
Inspired by PatchCore~\citep{roth2022patchcore}, this expert extracts dense patch tokens using DINOv2~\citep{oquab2024dinov2} at $518 \times 518$ resolution ($37 \times 37$ grid, $M = 1{,}369$ patches per image).
For each query patch $\patch^q_i$, we compute the cosine distance to its nearest neighbor in the reference patch bank $\mathcal{P}_{\text{ref}} = \bigcup_{I^r \in \refset} \{\patch^r_1, \ldots, \patch^r_M\}$:
\begin{equation}
    d_i = 1 - \max_{\patch \in \mathcal{P}_{\text{ref}}} \cos(\patch^q_i, \patch).
    \label{eq:patch_distance}
\end{equation}
The expert anomaly score aggregates the top-$1\%$ most anomalous distances:
\begin{equation}
    \ascore_{\text{patch}} = \frac{1}{|\mathcal{T}|} \sum_{i \in \mathcal{T}} d_i, \quad \mathcal{T} = \{i : d_i \geq d_{(0.99M)}\}.
    \label{eq:expert_score}
\end{equation}
The score is converted to a structured text report including the raw distance, global CLS-token similarity, and a qualitative assessment (``very similar to normal'' through ``strong anomaly signal'').

\paragraph{Visual retrieval expert.}
This module retrieves the $k$ most similar normal images from $\bank_\domcode$ using DINOv2 CLS-token cosine similarity and reports the retrieval scores, providing a global similarity context that complements the patch-level analysis.

\paragraph{Additional experts (extensible).}
The pool can include texture-statistics experts (Gram matrices, frequency-domain analysis), CLIP-based semantic similarity, or domain-specific modules (e.g., Hounsfield-unit windowing for CT images).
Each expert follows the same interface: accept image and references, return structured text.
The autonomous controller (\Cref{sec:controller}) decides which experts to invoke per query.

\subsection{Adversarial Debate Engine}
\label{sec:debate}

The core of AnomalyClaw is a structured debate between two VLM-based agents with opposing objectives (\Cref{fig:overview}).

\paragraph{Anomaly Proposer ($\proposer$).}
Given the query image, reference images, domain knowledge, and expert evidence, the Proposer generates a set of anomaly claims $\claim = \{c_1, \ldots, c_n\}$.
Each claim is a structured record:
\begin{equation}
    c_j = (\texttt{id}_j,\; \texttt{location}_j,\; \texttt{appearance}_j,\; \texttt{evidence}_j,\; \confidence_j,\; \texttt{severity}_j),
    \label{eq:claim}
\end{equation}
where $\texttt{location}_j$ specifies the spatial region (bounding box or quadrant), $\texttt{appearance}_j$ describes the visual observation, $\texttt{evidence}_j$ links the claim to specific expert findings, $\confidence_j \in [0,1]$ is the Proposer's confidence, and $\texttt{severity}_j$ is the assessed impact.
If no anomalies are found, the Proposer outputs a normal profile describing the expected appearance of the object.

\paragraph{Normality Advocate ($\advocate$).}
For each claim $c_j$, the Advocate generates a counter-argument:
\begin{equation}
    a_j = (\texttt{id}_j,\; \refuteconf_j,\; \texttt{counter\_evidence}_j,\; \texttt{likely\_cause}_j),
    \label{eq:refutation}
\end{equation}
where $\refuteconf_j \in [0,1]$ is the Advocate's confidence that the observation is \emph{not} anomalous, $\texttt{counter\_evidence}_j$ provides specific reasons (``lighting reflection matches reference angle,'' ``texture variation is within normal manufacturing tolerance''), and $\texttt{likely\_cause}_j$ attributes the observation to a benign source.
The Advocate's sole objective is to argue against anomaly claims, ensuring that only robust claims survive.

\paragraph{Claim resolution.}
After each debate round, claims are resolved by comparing the Proposer's confidence against the Advocate's refutation confidence:
\begin{equation}
    \text{status}(c_j) =
    \begin{cases}
        \texttt{Valid} & \text{if } \confidence_j \geq \tau_{\text{accept}} \text{ and } \refuteconf_j \leq \tau_{\text{reject}}, \\
        \texttt{Invalid} & \text{if } \refuteconf_j \geq \tau_{\text{accept}}, \\
        \texttt{TBD} & \text{otherwise},
    \end{cases}
    \label{eq:resolution}
\end{equation}
where $\tau_{\text{accept}} = 0.6$ and $\tau_{\text{reject}} = 0.4$ are resolution thresholds.
Claims with status \texttt{TBD} are carried into the next debate round for re-examination.

\paragraph{Multi-round refinement.}
When \texttt{TBD} claims remain after the first round, the Proposer and Advocate engage in additional rounds focused exclusively on unresolved claims.
The Proposer may revise confidence levels based on the Advocate's counter-arguments, and the Advocate may strengthen or weaken its challenges based on new evidence.
The debate terminates when all claims reach \texttt{Valid} or \texttt{Invalid} status, or when the maximum depth $D$ (controlled by the autonomous controller) is reached.

\begin{algorithm}[t]
\caption{AnomalyClaw adversarial debate}
\label{alg:debate}
\begin{algorithmic}[1]
\REQUIRE Query image $\query$, domain $\domcode$, normal bank $\bank_\domcode$, expert pool $\Phi$, max depth $D$
\ENSURE Verdict $v$, claim decisions
\STATE \textbf{// Expert Evidence Collection}
\STATE $\evidence \gets \controller.\text{SelectExperts}(\query, \domcode, \Phi)$ \hfill\COMMENT{Autonomous expert selection}
\STATE $\text{reports} \gets \{\expert(\query, \refset, \bank_\domcode) : \expert \in \evidence\}$
\STATE \textbf{// Debate Initialization (Round 0)}
\STATE $\claim \gets \proposer(\query, \refset, \text{reports}, \domcode)$ \hfill\COMMENT{Propose anomaly claims}
\STATE $\text{rebuttals} \gets \advocate(\query, \refset, \claim)$ \hfill\COMMENT{Challenge claims}
\STATE $\text{decisions} \gets \text{Resolve}(\claim, \text{rebuttals})$ \hfill\COMMENT{Eq.~\ref{eq:resolution}}
\STATE \textbf{// Multi-Round Refinement}
\FOR{$t = 1, \ldots, D-1$}
    \STATE $\text{TBD} \gets \{c_j : \text{decisions}[j] = \texttt{TBD}\}$
    \IF{$\text{TBD} = \emptyset$}
        \STATE \textbf{break} \hfill\COMMENT{All claims resolved}
    \ENDIF
    \STATE $\claim' \gets \proposer.\text{Refine}(\text{TBD}, \text{rebuttals})$
    \STATE $\text{rebuttals}' \gets \advocate.\text{Refine}(\text{TBD}, \claim')$
    \STATE $\text{decisions} \gets \text{Resolve}(\claim', \text{rebuttals}')$
\ENDFOR
\STATE \textbf{// Final Verdict}
\STATE $v \gets \controller.\text{Verdict}(\text{decisions})$
\RETURN $v$, decisions
\end{algorithmic}
\end{algorithm}

\paragraph{Why asymmetric debate avoids collective delusion.}
Symmetric multi-agent debate~\citep{li2025dmad,liang2026m3mad} assigns identical roles to all agents, which can lead to ``collective delusion''---agents reinforcing each other's errors through echo-chamber dynamics.
AnomalyClaw's asymmetric design avoids this by construction.
The Advocate's \emph{sole objective} is to argue against anomaly claims; it succeeds by finding convincing counter-evidence, not by agreeing with the Proposer.
This structural asymmetry ensures that false positives face genuine resistance, while true anomalies that survive the Advocate's challenges are accepted with higher confidence than single-pass judgments.

\subsection{Autonomous Controller}
\label{sec:controller}

The autonomous controller $\controller$ orchestrates the debate process, making three decisions:

\paragraph{Expert selection.}
Not all experts are informative for every query.
The controller first runs a lightweight assessment (e.g., global CLS-token similarity from the retrieval expert) and selects additional experts based on the initial signal:
\begin{itemize}[leftmargin=*,itemsep=1pt]
    \item If global similarity is high (query closely matches references), the patch-level expert is likely unnecessary for easy-normal cases.
    \item If global similarity is ambiguous, all available experts are invoked to provide maximum evidence.
    \item Domain-specific experts (e.g., CT windowing) are invoked only for relevant domains.
\end{itemize}
This adaptive selection reduces computational cost on easy queries while providing comprehensive evidence for difficult ones.

\paragraph{Debate depth control.}
The controller sets the maximum debate depth $D$ based on the initial assessment:
\begin{equation}
    D = \begin{cases}
        1 & \text{if initial confidence} \geq \tau_{\text{easy}} \text{ (clear case)}, \\
        D_{\max} & \text{otherwise (ambiguous case)},
    \end{cases}
    \label{eq:depth}
\end{equation}
where $\tau_{\text{easy}}$ is a confidence threshold (typically 0.8) and $D_{\max}$ is the maximum allowed depth (typically 2--3).
Easy cases---where the Proposer and Advocate quickly agree---terminate in a single round, saving the cost of additional VLM calls.

\paragraph{Final verdict.}
The controller produces the final anomaly verdict from the debate outcomes:
\begin{equation}
    v = \begin{cases}
        \texttt{anomalous} & \text{if any claim has status \texttt{Valid}}, \\
        \texttt{normal} & \text{if all claims are \texttt{Invalid} or no claims proposed}, \\
        \texttt{uncertain} & \text{if \texttt{TBD} claims remain at depth } D.
    \end{cases}
    \label{eq:verdict}
\end{equation}
Uncertain cases can be routed to human review or flagged for additional analysis, providing a natural interface for human-in-the-loop deployment.

\subsection{VLM-Agnostic Design}
\label{sec:vlm_agnostic}

AnomalyClaw is designed to work with any VLM that supports multi-image input and structured (JSON) output.
The Proposer and Advocate communicate through structured JSON schemas (\Cref{eq:claim,eq:refutation}), and expert evidence is provided as plain text.
This design avoids dependence on any specific model's capabilities or API format.
In our experiments (\Cref{sec:experiments}), we evaluate with \tbd{multiple VLM backends} and show consistent improvements across all of them.
